\subsection{Inserción de torres, marginal de dos casillas y menor lejano}\label{app:insercion}
Esta subsección enuncia y demuestra las tres estimaciones de torres citadas arriba: el teorema de
inserción del que proceden $\hat c_e$ y $r_e$, su corolario para dos casillas y la estabilidad (S) del
resto de inserción bajo un menor lejano. Aquí $M$ es una máscara con pesos de tamaño $k$ cuyas casillas
son las $e$ con $w_e>0$, $H=R\sqcup C$ es su conjunto de líneas, $a_e\in\mathbb R^H$ es el vector de
incidencia de $e$, $v=(\mathbf1_R,-\mathbf1_C)$ y $\Pi=vv^\top/(2k)$. El punto estacionario se escribe
$t_e:=p_e(1-p_e)=w_e\exp(a_e^\top\theta)$ con $\theta\in\mathbb R^H$, masas de línea uno y
$0<p_e<\tfrac12$. Además de $\tau_e=p_e(1-p_e)/(1-2p_e)$ usamos
\[
 \varsigma_e=\tau_e-p_e=\frac{p_e^2}{1-2p_e},\qquad m_a=\sum_{e\in a}\tau_e=1+\mu_a,\qquad
 \mu_a=\sum_{e\in a}\varsigma_e,
\]
y la función $c_2(u)=\log(1-2u)-2\log(1-u)$. Las letras $\varsigma,\mu,m$ sólo tienen este significado en
esta subsección. La hipótesis (H) con constantes $(C_p,\gamma_A)$, $0<\gamma_A\le1$, es $p_e\le C_p/k$
en toda casilla y $\lambda_{\min}(S_T+\Pi)\ge\gamma_A$. Sean
\begin{gather*}
 K_G=1+\frac{1+\frac76\sqrt{2C_p}}{\gamma_A},\qquad g_G=\tfrac76C_pK_G+\tfrac12K_G+\tfrac43C_p,\qquad
 C_\delta=\tfrac7{12}(4g_G+5C_p)\bigl(\tfrac73K_G+1\bigr),\\
 k_1=\max\Bigl\{41K_G^2C_p,\ 40g_G,\ \tfrac76(4g_G+5C_p)\bigl(\tfrac73K_G+1\bigr),\
            \frac{34K_GC_p}{\gamma_A}\Bigr\}.
\end{gather*}

\begin{teorema}[Inserción de torres]\label{teo:insercion-torres}
Sea $M$ una máscara que cumple \textup{(H)} con constantes $(C_p,\gamma_A)$ y sea $k\ge k_1$. Para toda
casilla $e$, con $\pi^{\rm der}_e=\partial_{\log w_e}L_T=p_e(1+\hat c_e)$ y
$r_e=L_T(M_e)-L_T(M)-\log(\pi^{\rm der}_e/w_e)$,
\[
 |\hat c_e|\le\frac{C_\delta}{k},\qquad |r_e|\le\frac{C_r}{k^2},
\]
donde $C_r$ sólo depende de $(C_p,\gamma_A)$. De modo equivalente, $w_ee^{L_T(M_e)-L_T(M)}=p_e(1+\hat
c_e)e^{r_e}$. En la clase $\mathcal K(w_0,w_1,M_0,\gamma)$ de la Sección~\ref{sec:torres}, la hipótesis
\textup{(H)} se cumple con $(C_p,\gamma_A)=(\tfrac32M_0,0{,}88\gamma)$ para $k\ge k_{\mathcal K}=
\max(31K_{\mathcal K}^2M_0,45K_{\mathcal K}M_0)$, $K_{\mathcal K}=1+(1+2{,}01\sqrt{M_0})/\gamma$. Por
tanto, las dos cotas valen en $\mathcal K$ para $k\ge\max\{k_{\mathcal K},k_1(\tfrac32M_0,0{,}88\gamma)\}$,
con constantes que sólo dependen de $(M_0,\gamma)$.
\end{teorema}
\begin{proof}
\emph{Paso 1. Cotas elementales.} Para $0<u\le\tfrac18$ se tiene $u\le\tau(u)\le\tfrac76u$ y
$u^2\le\varsigma(u)\le\tfrac43u^2$, donde $\tau(u)=u(1-u)/(1-2u)$ y $\varsigma=\tau-u$. Como
$k\ge k_1$ da $p_e\le C_p/k\le\tfrac18$, las masas de línea cumplen $1\le m_a\le\tfrac76$ y
$0\le\mu_a\le\tfrac43C_p/k$. Escribimos $a=\max_ep_e\le C_p/k$ y $\bar\tau=\max_e\tau_e$.

\emph{Paso 2. La matriz de Green.} Sea $G^0=(S_T+\Pi)^{-1}$. Entonces
\begin{equation}\label{eq:green-torres}
 \|G^0\|_\infty\le K_G,\qquad |G^0_{ab}-\delta_{ab}|\le\frac{g_G}{k},\qquad
 \sum_{e'}\tau_{e'}|\mathsf G(e,e')|\le\tfrac73K_G .
\end{equation}
En efecto, escribamos $S_T+\Pi=\Delta+W'$ con $\Delta=\operatorname{diag}(m_a)$. Si $(S_T+\Pi)x=b$,
entonces $\|x\|_2\le\sqrt{2k}\,\|b\|_\infty/\gamma_A$. La fila $a$ de $W'$ consta de las entradas
$\tau_e$, $e\in a$, y de la fila de $\Pi$, luego $\|W'_{a\cdot}\|_2\le(\bar\tau m_a)^{1/2}+(2k)^{-1/2}$,
y $m_a|x_a|\le|b_a|+\|W'_{a\cdot}\|_2\|x\|_2$. Con $k\bar\tau\le\tfrac76C_p$ resulta la primera cota. De
$\Delta G^0=I-W'G^0$ se obtiene $|G^0_{ab}-\delta_{ab}/m_a|\le\max_c|W'_{ac}|\sum_c|G^0_{cb}|/m_a
\le(\bar\tau+\tfrac1{2k})K_G$, y $|1/m_a-1|\le\mu_a$. Por último,
$\sum_{e'}\tau_{e'}|\mathsf G(e,e')|\le\sum_am_a|(G^0a_e)_a|\le\tfrac76\cdot2K_G$.

\emph{Paso 3. Función dual, identidades de Ward y $\hat c_e$.} Sea $u(t)=(1-\sqrt{1-4t})/2$ la raíz
pequeña de $u(1-u)=t$ y $\phi(t)=\log(1-u(t))+2u(t)$, de modo que $t\phi'(t)=u(t)$. Para una máscara
$N$ de tamaño $n$ pongamos
\[
 g_N(\theta)=\sum_e\phi\bigl(w_ee^{a_e^\top\theta}\bigr)-\sum_a\theta_a-2n .
\]
En $\{t_e(\theta)<\tfrac14\}$ es convexa, con $\nabla g_N=(\sum_{e\in a}p_e-1)_a$,
$\nabla^2g_N=S_T$ evaluado en $\tau(\theta)$, y $g_N(\theta+sv)=g_N(\theta)$. En un punto estacionario,
$-p_e\log(p_e/w_e)=p_e\log(1-p_e)-p_ea_e^\top\theta$, luego $H_T(p)=\sum_e\log(1-p_e)-\sum_a\theta_a
=g_N(\theta)=\min g_N$. La convexidad estricta en $v^\perp$ da la unicidad del punto estacionario.

Como $S_Tv=0$, se tiene $G^0v=v$ y $S_TG^0=G^0S_T=I-\Pi$. Puesto que $\Pi a_e=0$, para toda línea $a$ y
toda casilla $e$
\begin{equation}\label{eq:ward-torres}
 \sum_{e'\in a}\tau_{e'}\mathsf G(e',e)=\mathbf1_{\{e\in a\}},\qquad
 \sum_{e'}\tau_{e'}\mathsf G(e,e')=1 .
\end{equation}
Derivar $p_{e'}(1-p_{e'})=w_{e'}e^{a_{e'}^\top\theta}$ y las masas de línea respecto de $\log w_e$ da
$d\theta=-\tau_eG^0a_e\,d\log w_e$ en el calibre ortogonal a $v$, y
\[
 \frac{\partial p_{e'}}{\partial\log w_e}=\tau_e\delta_{ee'}-\tau_e\tau_{e'}\mathsf G(e,e'),
\]
cuya suma sobre cada línea se anula por~\eqref{eq:ward-torres}. Por el párrafo anterior, $\partial
H_T/\partial\log w_e=p_e$. Además $\partial\log\det(S_T+\Pi)=\sum_{e'}\dot\tau(p_{e'})\mathsf
G(e',e')\,\partial p_{e'}$ y $\partial\sum_{e'}c_2(p_{e'})=\sum_{e'}\dot c_2(p_{e'})\,\partial p_{e'}$; un
punto denota la derivada respecto de $u$, de modo que $\dot\tau=\kappa_2$ y $\dot c_2=-2\ell_2'$ en la
notación de la Proposición~\ref{prop:una-linea}.
Con $\varepsilon_{e'}=\dot c_2(p_{e'})+\dot\tau(p_{e'})\mathsf G(e',e')-2$, donde la constante $2$ se
puede restar porque las derivadas de $p$ suman cero, obtenemos
\begin{equation}\label{eq:chat-torres}
 \pi_e^{\rm der}=p_e+\tfrac12\tau_eX_e,\qquad \hat c_e=\frac{\tau_e}{2p_e}X_e,\qquad
 X_e=\sum_{e'}\tau_{e'}\varepsilon_{e'}\mathsf G(e,e')-\varepsilon_e .
\end{equation}
Es la expresión de $\hat c_c$ dada antes del Lema~\ref{lem:linea-inversa}. Multiplicar por una constante
los pesos de una línea no cambia $p$ y suma su logaritmo a $L_T$; por tanto
$\sum_{e\in a}p_e\hat c_e=0$ en toda línea. Por~\eqref{eq:green-torres}, $|\mathsf G(e,e)-2|\le4g_G/k
\le\tfrac1{10}$. Desarrollar $\dot c_2$ y $\dot\tau$ da $\varepsilon_e=(\mathsf G(e,e)-2)+2p_e+\tilde
\varepsilon_e$ con $|\tilde\varepsilon_e|\le12p_e^2+3p_e|\mathsf G(e,e)-2|$, luego
$|\varepsilon_e|\le(4g_G+5C_p)/k$. Entonces~\eqref{eq:green-torres} da
$|X_e|\le(4g_G+5C_p)(\tfrac73K_G+1)/k$, y $\tau_e/(2p_e)\le\tfrac7{12}$ da $|\hat c_e|\le
C_\delta/k$.

\emph{Paso 4. El hijo.} Fijemos $e_0=(r,c)$ y escribamos $p_0,\tau_0,\varsigma_0$ para sus valores. Sean
$J=\{r,c\}$, $I=H\setminus J$, $F$ el conjunto de las casillas de la fila $r$ o de la columna $c$,
$F^*=F\setminus\{e_0\}$, y para $e\in F^*$ sea $\lambda(e)\in I$ su línea no compartida con $e_0$. El
hijo $M'=M_{e_0}$ tiene las casillas ajenas a $F$, las líneas $I$ y $\Pi'=v_Iv_I^\top/(2k-2)$; las primas
denotan sus datos estacionarios, y $\Delta p_e=p'_e-p_e$. Restringir $\theta$ a $I$ deja en $M'$ el
defecto de masa $b=\nabla g_{M'}(\theta|_I)$, con $b_{\lambda(e)}=-p_e$ para $e\in F^*$ y $b_a=0$ en las
demás líneas; $b\perp v_I$. Sea $S'(\tau|)$ el operador del hijo con los pesos del padre y
$\xi_N=-(S'(\tau|)+\Pi')^{-1}b$ el paso de Newton. Para $x\perp v_I$, extendido por cero a $\tilde
x\perp v$, se tiene $x^\top S'(\tau|)x=\tilde x^\top S_T\tilde x-\sum_{e\in F^*}\tau_ex_{\lambda(e)}^2
\ge(\gamma_A-\bar\tau)|x|^2$, y las masas de $S'(\tau|)$ son al menos $1-\bar\tau$. El Paso~2 aplicado a
$S'(\tau|)$ da $\|(S'(\tau|)+\Pi')^{-1}\|_\infty\le1{,}1K_G$. En la bola de radio $\tfrac1{16}$ la
hessiana $\nabla^2g_{M'}$ es de Lipschitz en la norma del supremo con constante a lo sumo $9$, porque
$dp/d\sigma=\tau$ y $d\tau/d\sigma\le1{,}8\tau$ a lo largo de $t\mapsto te^\sigma$. Como $41K_G^2a\le1$,
el teorema de Newton--Kantorovich en $v_I^\perp$ con la norma del supremo da un único punto estacionario
$\theta'=\theta|_I+\xi$ del hijo, $\xi\perp v_I$, con
\begin{equation}\label{eq:newton-hijo}
 \|\xi\|_\infty\le3K_Ga,\qquad \|\xi-\xi_N\|_\infty\le61K_G^3a^2 .
\end{equation}
Los desplazamientos $s_e=a_e^\top\xi$ cumplen $|s_e|\le6K_Ga\le\tfrac3{41}$, luego $p'_e\le1{,}09p_e$ y
$\tau'_e\ge0{,}9\tau_e$. Por tanto el hijo cumple \textup{(H)} en tamaño $k-1$ con constantes
$(\tfrac54C_p,\tfrac34\gamma_A)$.

\emph{Paso 5. Descomposición exacta de $r_{e_0}$.} Por el Paso~3 para $M$ y para $M'$,
$g_M(\theta)-g_{M'}(\theta|_I)=\sum_{e\in F}\phi(t_e)-\theta_r-\theta_c-2$, con
$\theta_r+\theta_c=\log(p_0(1-p_0)/w_{e_0})$, $\sum_F\phi(t_e)=\sum_F\log(1-p_e)+2(2-p_0)$ y
$\sum_{F^*}p_e=2-2p_0$. Por tanto
\begin{equation}\label{eq:accion-hijo}
 H_T(M')-H_T(M)-\log\frac{p_0}{w_{e_0}}=\sum_{e\in F^*}\bigl[-\log(1-p_e)-p_e\bigr]-\mathcal D,\qquad
 \mathcal D=g_{M'}(\theta|_I)-g_{M'}(\theta')\ge0 .
\end{equation}
Para el prefactor, sean $T_{JJ}=\bigl(\begin{smallmatrix}m_r&\tau_0\\ \tau_0&m_c\end{smallmatrix}\bigr)$,
$T_{IJ}$ el bloque de $S_T$, $D_F=\operatorname{diag}(\tau_e)_{a=\lambda(e)}$,
$Q_J=T_{IJ}T_{JJ}^{-1}T_{JI}$ y
\[
 \mathsf S=S'(\tau|)+D_F-Q_J,\qquad W=\mathsf S-S'(\tau') .
\]
Así $\mathsf S$ es el complemento de Schur de $S_T$ en $I$. Para una matriz semidefinida positiva con
núcleo generado por $v$, el menor principal que omite la línea $a$ vale el pseudodeterminante por
$v_a^2/|v|^2$. Aplicado a $S_T$ y a $\mathsf S$, con $|v|^2=2k$ y $|v_I|^2=2k-2$, da
$\mathsf Sv_I=Wv_I=0$ y $\operatorname{pdet}S_T=\tfrac k{k-1}\det T_{JJ}\operatorname{pdet}\mathsf S$.
Resolver $S_T(y,z)=(0,x)$ para $x\perp v_I$ da $\mathsf S^+=\Pi_I^\perp G^0_{II}\Pi_I^\perp$, luego
$\|\mathsf S^+\|_2\le\gamma_A^{-1}$ y $\operatorname{tr}(\mathsf S^+W)=\operatorname{tr}(G^0_{II}W)$.
El término $\tfrac12\log(4\pi k)$ de $L_T$ aporta $\tfrac12\log\tfrac{k-1}k$, que cancela el factor de
tamaño. Con $\ell=L_T-H_T$,
\[
 \ell(M')-\ell(M)=\tfrac12\log\det T_{JJ}+\tfrac12\sum_{e\in F}c_2(p_e)
   -\tfrac12\sum_{e\in M'}[c_2(p'_e)-c_2(p_e)]-\tfrac12\log\det(I-\mathsf S^+W).
\]
Al sumar~\eqref{eq:accion-hijo} obtenemos la identidad exacta
\begin{equation}\label{eq:re-exacto}
 r_{e_0}=-\mathcal D+L_1+L_2+L_3-\log(1+\hat c_{e_0}),
\end{equation}
donde $L_1=\sum_{F^*}[-\log(1-p_e)-p_e]+\tfrac12\log\det T_{JJ}+\tfrac12\sum_Fc_2(p_e)$,
$L_2=-\tfrac12\sum_{M'}[c_2(p'_e)-c_2(p_e)]$ y $L_3=-\tfrac12\log\det(I-\mathsf S^+W)$.

Sea $\beta=G^0a_{e_0}$, de modo que $S_T\beta=a_{e_0}$, $v^\top\beta=0$ y $\mathsf G_0:=\mathsf
G(e_0,e_0)=\beta_r+\beta_c$. Las filas $a\in I$ de $S_T\beta=a_{e_0}$ muestran que
$\rho:=S'(\tau|)\beta_I-b$ es ortogonal a $v_I$, se anula fuera de las líneas $\lambda(e)$, $e\in F^*$,
y
\begin{equation}\label{eq:rho-hijo}
 \rho_{\lambda(e)}=-\varsigma_e-\tau_e(\mathsf G(e,e_0)-1)\quad(e\in F^*),\qquad
 \zeta:=(S'(\tau|)+\Pi')^{-1}\rho+(\xi-\xi_N).
\end{equation}
Como $S'(\tau|)(\xi_N+\beta_I)=\rho$, el vector $\zeta-(\xi+\beta_I)$ es un múltiplo de $v_I$, y
$a_e^\top\zeta=s_e+\mathsf G(e,e_0)$ para toda casilla $e$ de $M'$. Por~\eqref{eq:green-torres}
y~\eqref{eq:newton-hijo}: $|\beta_a|\le2g_G/k$ para $a\in I$, $|\beta_j-1|\le2g_G/k$ para $j\in J$,
$\sum_a|\beta_a|\le2K_G$; $|\mathsf G(e,e_0)-1|\le4g_G/k$ en $F^*$ y $|\mathsf G(e,e_0)|\le4g_G/k$ en
$M'$; y
\begin{equation}\label{eq:cotas-hijo}
 \|\rho\|_\infty\le\tfrac43a^2+\frac{4\bar\tau g_G}k,\qquad
 \|\rho\|_1\le\tfrac83a+\frac{28g_G}{3k},\qquad
 \|\zeta\|_\infty\le\tfrac32K_G\|\rho\|_\infty+61K_G^3a^2 .
\end{equation}

\emph{Paso 6. Cancelación de primer orden.} Sean
\begin{gather*}
 S_1=\sum_{M'}p_e\tau_e\mathsf G(e,e_0),\qquad
 S_2=\sum_{M'}\tau_e\mathsf G(e,e_0)(\mathsf G(e,e)-2),\\
 S_3=\sum_{F^*}\tau_e(\mathsf G(e,e)-2),\qquad S_4=\sum_{F^*}\tau_ep_e .
\end{gather*}
Las identidades siguientes \emph{definen} los restos $E_\bullet$:
\begin{align*}
 \mathcal D&=\tfrac12(m_r\beta_r+m_c\beta_c)-1+\tfrac12\tau_0\mathsf G_0+E_{\mathcal D}, &
 L_1&=\tfrac12(\mu_r+\mu_c)+E_{L1},\\
 L_2&=-S_1+E_{L2}, &
 L_3&=\tfrac12\operatorname{tr}(G^0_{II}W)+E_{L3},\\
 -\log(1+\hat c_{e_0})&=-\hat c_{e_0}+E_{\delta1}, &
 \hat c_{e_0}&=\tfrac12X_{e_0}+E_{\delta2},
\end{align*}
junto con $\operatorname{tr}(G^0_{II}W)=\mathrm{(i)}+\mathrm{(ii)}-\mathrm{(iii)}$, donde
$\mathrm{(i)}=-\sum_{M'}(\tau'_e-\tau_e)\mathsf G(e,e)=-2(1-p_0)+S_2+4S_1+E_{\rm(i)}$,
$\mathrm{(ii)}=\sum_{F^*}\tau_eG^0_{\lambda(e)\lambda(e)}$ y
$\mathrm{(iii)}=\operatorname{tr}(G^0_{II}Q_J)=\tfrac2k-2+m_rG^0_{rr}+m_cG^0_{cc}+2\tau_0G^0_{rc}$, y
con $X_{e_0}=S_3+2S_4+S_2+2S_1-(\mathsf G_0-2)-2p_0+E_X$. La primera identidad es el desarrollo de
Taylor de $g_{M'}$ en $\theta|_I$, es decir, $\mathcal D=-b^\top\xi-\tfrac12\xi^\top S'(\tau|)\xi-R_3$
con $R_3=\sum_{M'}[\phi(t_ee^{s_e})-\phi(t_e)-p_es_e-\tfrac12\tau_es_e^2]$. Sustituyendo
$\xi\equiv\zeta-\beta_I$ y $S'(\tau|)\beta_I=b+\rho$, y usando las filas $r,c$ de
$S_T\beta=a_{e_0}$ en la forma $\sum_{F^*}\tau_e\beta_{\lambda(e)}=2-m_r\beta_r-m_c\beta_c-\tau_0\mathsf
G_0$, resulta
\[
 E_{\mathcal D}=\tfrac12\sum_{F^*}\varsigma_e\beta_{\lambda(e)}-\tfrac12\beta_I^\top\rho+\zeta^\top\rho
   -\tfrac12\zeta^\top S'(\tau|)\zeta-R_3 .
\]
El desarrollo de $\mathrm{(i)}$ usa la conservación exacta de la masa $\sum_{M'}\Delta p_e=1-p_0$,
$\tau=p+\varsigma$ y $s_e=-\mathsf G(e,e_0)+a_e^\top\zeta$; el de $X_{e_0}$
usa~\eqref{eq:chat-torres} y $\varepsilon_e=(\mathsf G(e,e)-2)+2p_e+\tilde\varepsilon_e$. El valor de
$\mathrm{(iii)}$ se sigue de los bloques $(J,I)$ y $(J,J)$ de $S_TG^0=G^0S_T=I-\Pi$.

En la suma de los términos principales de~\eqref{eq:re-exacto}, $S_1$ tiene coeficiente $-1+2-1=0$ y
$S_2$ tiene $\tfrac12-\tfrac12=0$: todo lo que no toca la fila $r$ ni la columna $c$ se cancela. Para
$e\in F^*$ con línea compartida $j$ y $\lambda=\lambda(e)$ se tiene $\mathsf
G(e,e)=G^0_{jj}+G^0_{\lambda\lambda}+2G^0_{j\lambda}$, y la entrada $(r,r)$ de $S_TG^0=I-\Pi$ da
$\sum_{e\in r\setminus e_0}\tau_eG^0_{c(e)r}=1-\tfrac1{2k}-m_rG^0_{rr}-\tau_0G^0_{cr}$, e igualmente para
$c$. Con estas relaciones los términos de orden $1/k$ se cancelan exactamente, y
\begin{align}
 r_{e_0}&=E_{\rm main}-E_{\mathcal D}+E_{L1}+E_{L2}+\tfrac12E_{\rm(i)}+E_{L3}
         +E_{\delta1}-E_{\delta2}-\tfrac12E_X,\label{eq:re-restos}\\
 E_{\rm main}&=\sum_{e\in r\setminus e_0}\frac{p_e^3}{1-2p_e}+\sum_{e\in c\setminus e_0}\frac{p_e^3}{1-2p_e}
   -\tfrac12(G^0_{rr}-1)\mu_r-\tfrac12(G^0_{cc}-1)\mu_c-\tfrac12G^0_{rc}(\mu_r+\mu_c).\notag
\end{align}

\emph{Paso 7. Cotas.} Cada resto de~\eqref{eq:re-restos} es a lo sumo $k^{-2}$ por un polinomio en
$(C_p,K_G,g_G,\gamma_A^{-1})$. Para $E_{\rm main}$ esto se sigue de $\sum_{e\in r}p_e^3/(1-2p_e)\le
\tfrac43a^2$, $|G^0-I|\le g_G/k$ y $\mu_a\le\tfrac43a$. Para $E_{\mathcal D}$ usamos
$\sum_{F^*}\varsigma_e\le\tfrac83a$, las cotas de $\beta$ y~\eqref{eq:cotas-hijo}, $\zeta^\top
S'(\tau|)\zeta=\sum_{M'}\tau_e(a_e^\top\zeta)^2$ con $\sum_e\tau_e\le\tfrac76k$, y
$|\tfrac{d^3}{ds^3}\phi(te^s)|=|\tfrac{d\tau}{ds}|\le1{,}8\tau$, que da $|R_3|\le0{,}35k(6K_Ga)^3$. Para
$E_{L1}$: $|c_2(u)+u^2|\le3{,}6u^3$, $0\le-\log(1-u)-u-\tfrac12u^2\le0{,}39u^3$,
$|\log m_a-\mu_a|\le\tfrac12\mu_a^2$ y $\tau_0^2\le(\tfrac76a)^2$. Para $E_{L2}$: $|\dot c_2(u)/2+u|\le
4{,}2u^2$, $|\Delta p_e-\tau_es_e|\le0{,}9\tau_es_e^2$ y $|\ddot c_2|\le5{,}3$. Para $E_{\rm(i)}$:
$\Delta\varsigma=\dot\varsigma\Delta p+\tfrac12\ddot\varsigma\Delta p^2$ y $|\dot\varsigma(p)\mathsf
G(e,e)-4p|\le22{,}5p^2+8pg_G/k$. Para $E_{L3}$: con $Y=(\mathsf S^+)^{1/2}W(\mathsf S^+)^{1/2}$ se tiene
$\|Y\|\le\|W\|_\infty/\gamma_A\le\tfrac12$, $\|Y\|_F\le\|W\|_F/\gamma_A$ y
$|\log\det(I-Y)+\operatorname{tr}Y|\le\|Y\|_F^2$. Por la conservación de la masa en cada línea $a\in I$,
las entradas de $W$ son $W_{aa}=\varsigma_{e_a}-\sum_{e\in a\cap M'}\Delta\varsigma_e-(Q_J)_{aa}$, con
$e_a$ la casilla de $F^*$ en $a$, y $W_{ab}=-(\tau'_e-\tau_e)-(Q_J)_{ab}$ para la casilla $e$ de las
líneas $a,b$; por tanto $\max|W_{ab}|=O(k^{-2})$, $\|W\|_\infty=O(k^{-1})$ y $\|W\|_F^2=O(k^{-2})$. Por
último, $|E_{\delta1}|\le\hat c_{e_0}^2$, $|E_{\delta2}|\le\tfrac23a|X_{e_0}|$ porque
$\tau_0/p_0-1\le\tfrac43a$, y $E_X$ reúne $\tau_0\varepsilon_0\mathsf G_0$, las diferencias $\mathsf
G(e,e_0)-1$ en $F^*$ y los términos $\tilde\varepsilon$, todos controlados por el Paso~3. Sumar las
cotas en $k=k_1$ y usar que cada una es $k^{-2}$ o $k^{-3}$ por una constante da $|r_{e_0}|\le C_r/k^2$
para todo $k\ge k_1$. El enunciado para $\mathcal K$ se sigue de las propiedades de la clase
estacionaria recordadas en la Sección~\ref{sec:torres}.
\end{proof}

